% Quran-Tajzee: IEEE conference paper (4 pages max). Build: make paper (XeLaTeX).
\documentclass[conference]{IEEEtran}
\IEEEoverridecommandlockouts
\usepackage{fontspec}
\usepackage{amsmath,amssymb}
\usepackage{booktabs}
\usepackage{cite}
\usepackage{pgfplots}
\pgfplotsset{compat=1.18}
\usepackage{url}
\usepackage{polyglossia}
\setmainlanguage{english}
\setotherlanguage{arabic}
\newfontfamily\arabicfont[Script=Arabic,Scale=1.05]{Amiri}
\newcommand{\ar}[1]{\mbox{\textarabic{#1}}}
% Generated by scripts/paper_numbers.py; do not edit by hand.
\newcommand{\Ayat}{6,236}
\newcommand{\Phrases}{1,171}
\newcommand{\PhrasesReading}{1,094}
\newcommand{\PhraseMean}{5.3}
\newcommand{\PearsonAya}{0.990}
\newcommand{\PearsonPhrase}{0.994}
\newcommand{\RecitationTotal}{398,670}
\newcommand{\Letters}{325,386}
\newcommand{\JuzOff}{2.4}
\newcommand{\RubOff}{15}
\newcommand{\RubFloor}{1.5}
\newcommand{\RubExcess}{10}
\newcommand{\RubLetterSpread}{2.2}
\newcommand{\TradCuts}{90}
\newcommand{\TradInnerCuts}{79}
\newcommand{\TradJuzCuts}{11}
\newcommand{\CutDepth}{50}
\newcommand{\TradLeaning}{21}
\newcommand{\TradHizbWorst}{13.9}
\newcommand{\TradHizbCV}{6.9}
\newcommand{\TradRubWorst}{38.3}
\newcommand{\TradRubCV}{11.7}
\newcommand{\TradRubRatio}{2.19}
\newcommand{\TradHizbRatio}{1.36}
\newcommand{\NewHizbWorst}{3.6}
\newcommand{\NewHizbCV}{1.6}
\newcommand{\NewRubWorst}{30.7}
\newcommand{\NewRubCV}{6.9}
\newcommand{\NewRubRatio}{1.77}
\newcommand{\NewHizbRatio}{1.17}
\newcommand{\TradJuzWorst}{6.2}
\newcommand{\OptJuzWorst}{5.0}
\newcommand{\NewInnerCuts}{0}
\newcommand{\HizbCandidates}{61}
\newcommand{\HizbSplits}{7}
\newcommand{\JuzSplits}{11}
\newcommand{\RubRoundBoundaries}{137}
\newcommand{\RubRoundCandidates}{464}
\newcommand{\RubRoundSplits}{59}
\newcommand{\RubRoundKept}{78}
\newcommand{\ReviewedRub}{(59,6.90)}
\newcommand{\ReviewedRubWorst}{(59,30.70)}
\newcommand{\CeilRubWorst}{20.8}
\newcommand{\CeilRubCV}{3.4}
\newcommand{\CurveCV}{(0,9.63)(28,6.22)(53,4.83)(63,4.48)(102,3.62)(114,3.48)(131,3.40)}
\newcommand{\CurveWorst}{(0,30.69)(28,20.78)(53,20.78)(63,20.78)(102,20.78)(114,20.78)(131,20.78)}
\newcommand{\CurveRows}{0 & $\pm$30.7\% & 9.6\% \\ 28 & $\pm$20.8\% & 6.2\% \\ 53 & $\pm$20.8\% & 4.8\% \\ 63 & $\pm$20.8\% & 4.5\% \\ 102 & $\pm$20.8\% & 3.6\% \\ 114 & $\pm$20.8\% & 3.5\% \\ 131 & $\pm$20.8\% & 3.4\%}
\newcommand{\KneeSplits}{28}
\newcommand{\KneeWorst}{20.8}
\newcommand{\KneeCV}{6.2}


\begin{document}

\title{Quran-Tajzee: Balanced Quran Divisions\\ that End Where a Thought Ends}

\author{\IEEEauthorblockN{Tarek Eldeeb}
\IEEEauthorblockA{Giza, Egypt\\
tarekeldeeb@gmail.com}}

\maketitle

\begin{abstract}
The Madina Mushaf divides the Quran into 30 juz, 60 hizb and 240 quarters (rub\textsuperscript{c}),
taken from al-Ṣafāqusī's \emph{Ghayth al-Naf\textsuperscript{c}}. His stated rule is that each
quarter ends at a complete stop, at a position that is agreed or well known; the quarters are
not computed. We measure the transmitted divisions on the Tanzil text with two length models,
Hafs recitation in harakat with madd lengths and DigitalKhatt print width, which correlate at
$r=\PearsonAya$ per aya. The juz lie within \JuzOff\% of equal-letter positions, but the quarters
lie \RubExcess$\times$ further from equal positions than aya ends force and vary
\RubLetterSpread$\times$ in length; and \TradCuts{} of 239 boundaries fall inside a subject
phrase. Keeping the traditional juz, we re-cut hizb and quarters by an exact dynamic program
restricted to the starts of \Phrases{} reviewed subject phrases of 1--10 ayat. The worst hizb
drops from $\pm$\TradHizbWorst\% to $\pm$\NewHizbWorst\% of half a juz (CV \TradHizbCV\%
$\rightarrow$ \NewHizbCV\%), and no new boundary falls inside a phrase. Quarter balance is
limited by phrase granularity; we show that about a hundred targeted, text-checked phrase
splits bring the quarter CV from \TradRubCV\% to near the aya-level floor of \CeilRubCV\%.
\end{abstract}

\begin{IEEEkeywords}
Quran, text segmentation, partitioning, dynamic programming, typesetting, recitation, tajweed
\end{IEEEkeywords}

\section{Introduction}
Readers, teachers and printers organise the Quran by its divisions: a juz a day completes it in
a month, a hizb or a quarter is a daily portion or a unit of revision, and printed Mushafs mark
every quarter. Those divisions are of scholarly judgement (\emph{ijtihādī}), not revelation, which
is why mushafs and regions differ on them~\cite{hamad}. The Madina Mushaf takes its 30 juz, 60
hizb and their quarters from \emph{Ghayth al-Naf\textsuperscript{c}} by al-Ṣafāqusī (d.\ 1118\,AH)~\cite{ghayth,madina}.

Two properties are wanted of a division: units of similar length, and units that are complete in
meaning. The juz satisfy the first; the quarters satisfy neither fully (Section~\ref{sec:trad}).
This paper (i) measures the traditional divisions against their own stated method,
(ii) defines a length model for recitation and print, (iii) solves the hizb and quarters inside
the traditional juz exactly, with every boundary at the start of a subject phrase, and (iv) shows
how targeted phrase splits, each checked against the text, reach the balance that aya
granularity allows. Data, code and every number in this paper are in the open repository~\cite{repo}.

\section{The traditional method and its outcome}\label{sec:trad}
In the terminology section of his introduction al-Ṣafāqusī states his method~\cite{ghayth}: he
does not compute boundaries but follows only the agreed or well-known ones,
``\ar{ولا أمشي إلا على المتفق عليه أو المشهور}''; he treats each quarter as a unit,
``\ar{وأذكر حكم كل ربع بانفراده}''; and he closes each on a proper stop, naming its last word, the
kind of stop it allows and whether it is an aya end. At 2:25 he writes
``\ar{خالدون تام في أعلى درجاته وفاصلة ومنتهى الربع بإجماع}'': a complete stop of the highest degree,
an aya end, the end of the quarter by consensus. Where traditions differ he lists the positions
of the Maghāriba and the Mashāriqa and states a preference; the two cases we checked both
match the Madina Mushaf (45:12, 82:1).

\textbf{Drawn at stops, not by count.} The early divisions were made by counting letters~\cite{masahif,bayan}.
Counting the \Letters{} written letters of the text (basmala excluded), the traditional juz start on
average \JuzOff\% of a juz from equal-letter positions. Inside each juz, however, the quarters
start \RubOff\% of a quarter from equal positions, while the nearest aya end is only \RubFloor\%
away: \RubExcess$\times$ further than aya granularity forces. Quarter lengths vary
\RubLetterSpread$\times$ in letters. The quarters were placed at the ends of thoughts, as the method
says, not by count.

\textbf{But not always at the ends of thoughts.} We segmented the text into subject phrases
(Section~\ref{sec:data}). Against that reading, \TradCuts{} of the 239 traditional boundaries,
\TradInnerCuts{} of the 210 hizb and quarter boundaries inside the juz, fall inside a phrase,
at a median \CutDepth\% of the way through it. Some open mid-dialogue or mid-story: juz~16 at
``\ar{قال ألم أقل لك}'' (18:75), juz~27 at ``\ar{قال فما خطبكم أيها المرسلون}'' (51:31), a hizb at
``\ar{فنبذناه بالعراء}'' (37:145). The quarters thus pay the price of uneven length without
fully keeping the rule they were drawn by. This is the problem we address.

\section{Data and length model}\label{sec:data}
\textbf{Text.} Tanzil Uthmani 1.1, \Ayat{} ayat in the Kufan count~\cite{tanzil}.

\textbf{Subject phrases.} The text was divided into runs of 1--10 consecutive ayat on one subject,
drafted by reading the text and reviewed by hand: \PhrasesReading{} phrases as a reading of the
text, \Phrases{} after the review round of Section~\ref{sec:review}. A phrase start is our proxy
for a complete stop: the next aya does not continue the previous subject.

\textbf{Recitation size.} Per aya, the phonemes of the Hafs recitation from
\emph{tajweed-embeddings}~\cite{tajemb}: a madd letter counts its length in harakat (2, 4 or 6,
with \emph{tawassuṭ} for the optional lengths), a silent letter 0, every other letter 1. Each aya
is recited continuously and stopped at its end. A surah opening adds the basmala (19 harakat),
except At-Tawba. The whole text is \RecitationTotal{} harakat.

\textbf{Print size.} Per aya, the width of the aya set as a single line in the DigitalKhatt
Madina font~\cite{digitalkhatt}, shaped with HarfBuzz by Quran-Nazm~\cite{nazm}, with its aya
marker and one space.

The two sizes correlate at $r=\PearsonAya$ per aya and $r=\PearsonPhrase$ per phrase. We mix them
50/50 as shares of their totals, $c_a = \tfrac12 r_a/R + \tfrac12 p_a/P$.

\section{Method}
\textbf{Keep the juz.} An exact solver over the phrases (penalising mid-surah starts) can make the
30 juz more even, but only from $\pm$\TradJuzWorst\% to $\pm$\OptJuzWorst\%. The traditional juz are
kept; the imbalance is in the finer divisions.

\textbf{Nested exact cuts.} Each juz is cut into two hizb, then each hizb into four quarters, the
way the divisions nest. A container $G$ of ayat cut into $k$ parts has target $T=|G|/k$ and the
admissible cuts are phrase starts $\mathcal{S}$. Lexicographically we
\begin{equation}
\min_{\text{cuts}\subset\mathcal{S}}\ \max_q \frac{|S_q - T|}{T},\quad\text{then}\quad
\min \sum_q (S_q-T)^2 ,
\end{equation}
both by the dynamic program $F_k(i)=\min_{j<i}\,[F_{k-1}(j)\oplus g(C_i-C_j)]$ over prefix sums $C$,
with $\oplus=\max$ or $+$~\cite{bellman}. Contiguous partitioning has optimal substructure, so the
solution is the global optimum; tests check it against exhaustive search. A \emph{penalized}
variant also admits any aya end at a cost $m^2$ per mid-phrase start; we use it only to
\emph{propose} phrase splits.

\textbf{Review process.}\label{sec:review} Balance is limited by how fine the phrases are, so
phrases are refined where the divisions need it, one container at a time. (1)~For every aya end
in the container, the balance is computed if the next unit started there. (2)~The candidates are
the aya ends closer to balance than the current phrase-start solution; all lie inside a phrase.
(3)~Each candidate is read in context and accepted only at a complete stop: the aya before it
closes its sentence and subject, and the aya after it does not lean on it. A reply (\ar{قال}),
a pointer back (\ar{أولئك}, \ar{الذين}), a \ar{ف}/\ar{ثم} continuation, or a description of what
precedes it disqualifies a candidate. (4)~Among accepted candidates the most even wins; if none
qualifies, the current start is kept, however uneven. (5)~An accepted candidate splits its phrase,
with a topic for each part, so the rule ``every boundary opens a phrase'' is never relaxed.
Every candidate, its balance, the decision and the reason are recorded in a table that a human
reviewer approves or overturns; the phrase list is rebuilt from the reading and the accepted rows.

\emph{Accepted (juz~10).} The second hizb started at 9:29, $-3.7\%$ off. Candidate 9:30, $-1.5\%$,
lies inside the phrase 9:29--35. The aya before it ends ``\ar{وهم صاغرون}'', a complete stop
closing the command to fight; 9:30 ``\ar{وقالت اليهود عزير ابن الله}'' opens a new subject, the
claims of the People of the Book. The phrase is split into 9:29 and 9:30--35.

\emph{Rejected (juz~6).} The second hizb starts at 5:27, $+3.6\%$. Candidate 5:24, $+0.2\%$, is
almost exactly at the middle, but ``\ar{قالوا يا موسى إنا لن ندخلها}'' is the people's reply
inside the dialogue of 5:20--26; every closer point is a turn of that dialogue. The start is kept
at 5:27, the story of the two sons of Ādam, at the cost of 3.4 points of balance.

\section{Results}
\textbf{Hizb.} One review round checked all 30 second-hizb starts: 10 were already at the best
point, \HizbSplits{} phrases were split at complete stops (e.g.\ 9:30 ``\ar{وقالت اليهود}'',
10:28 ``\ar{ويوم نحشرهم}''), and 13 were kept because every closer point lies inside a story or
dialogue (e.g.\ juz~6, +3.6\%, inside the dialogue of Mūsā and his people). The \JuzSplits{}
phrases cut by traditional juz starts were split too. Table~\ref{tab:bal} gives the balance.

\begin{table}[t]
\caption{Traditional and new divisions (50/50 size mix). Hizb against \(\tfrac12\) of their juz,
quarters against \(\tfrac14\) of their hizb; ratios over the whole Quran.}
\label{tab:bal}
\centering
\begin{tabular}{@{}lrr@{}}
\toprule
& Traditional & Quran-Tajzee\\
\midrule
Worst hizb & $\pm$\TradHizbWorst\% & $\pm$\NewHizbWorst\% \\
Hizb CV & \TradHizbCV\% & \NewHizbCV\% \\
Longest/shortest hizb & \TradHizbRatio$\times$ & \NewHizbRatio$\times$ \\
Worst quarter & $\pm$\TradRubWorst\% & $\pm$\NewRubWorst\% \\
Quarter CV & \TradRubCV\% & \NewRubCV\% \\
Longest/shortest quarter & \TradRubRatio$\times$ & \NewRubRatio$\times$ \\
Starts inside a phrase (of 210) & \TradInnerCuts & \NewInnerCuts \\
\bottomrule
\end{tabular}
\end{table}

\textbf{Quarters.} With phrase starts only, a hizb holds about 18 phrases, 4.6 per quarter, and
one long phrase can be a fifth of a quarter; the worst quarter improves only to
$\pm$\NewRubWorst\%. Fig.~\ref{fig:curve} shows what targeted splits buy: the penalized solver
proposes the splits, the phrase-only solver is re-run with them accepted. About a hundred splits
bring the quarter CV close to the floor of \CeilRubCV\% reached when every aya end is allowed;
the remaining worst case, $\pm$\CeilRubWorst\%, is set by single long ayat (2:282 alone is 40\% of
a quarter), and going lower would require cutting inside an aya, which no traditional division
does. Uniform re-segmentation is far less efficient: splitting every phrase to $\le$3 ayat adds
about 1,350 boundaries for a CV of 4.5\%.

\begin{figure}[t]
\centering
\begin{tikzpicture}
\begin{axis}[width=\columnwidth, height=4.2cm, xlabel={phrase splits accepted},
  ylabel={quarter deviation (\%)}, ymin=0, ymax=26, xmin=0, xmax=150, grid=major,
  legend style={font=\footnotesize, at={(0.98,0.97)}, anchor=north east},
  tick label style={font=\footnotesize}, label style={font=\footnotesize}]
\addplot+[mark=*, thick] coordinates {\CurveWorst};
\addplot+[mark=square*, thick] coordinates {\CurveCV};
\addplot[dashed, domain=0:150] {\CeilRubCV};
\legend{worst quarter, CV, CV floor (any aya end)}
\end{axis}
\end{tikzpicture}
\caption{Quarter balance against the number of targeted phrase splits.}
\label{fig:curve}
\end{figure}

\section{Discussion and limitations}
\textbf{The phrase test is ours.} That no new boundary falls inside a phrase holds by
construction; the phrases are one reading of the text, under review. A wording check independent
of it, a first word that leans on the previous aya (a reply such as ``\ar{قال}'' or a linking
``\ar{ف}'', ``\ar{ثم}''), flags \TradLeaning{} traditional starts and a similar number of new ones,
most of them a resumptive ``\ar{ف}'' at a new subject (``\ar{فاصدع بما تؤمر}''), but a few real: the
10-aya phrase limit splits long stories at turns of dialogue. An independent judge is the pause
grades of the classical \emph{waqf} works, al-Dānī's \emph{al-Muktafā} and al-Ushmūnī's
\emph{Manār al-Hudā}~\cite{muktafa,manar}, at every candidate aya end.

\textbf{Counting schools.} Boundaries are Kufan aya ends; al-Ṣafāqusī also requires an aya end
agreed by all counts (\ar{فاصلة بلا خلاف}), not yet checked. \textbf{Weights.} Recitation and
print agree so closely that the 50/50 mix is not sensitive; header lines of surahs, if counted in
print, lengthen juz~30.

\section{Conclusion}
Al-Ṣafāqusī's rule, that a quarter ends where a thought ends, is sound, but the transmitted
quarters keep it only partly and vary more than twofold in length. Applying the rule to every
boundary and choosing among the admissible stops by measured length gives hizb within
$\pm$\NewHizbWorst\% of half their juz, none inside a phrase. For the quarters, the next review round
will split about a hundred phrases at checked stops, which brings their spread close to what
aya granularity allows.

\begin{thebibliography}{99}
\bibitem{ghayth} ʿA. al-Nūrī al-Ṣafāqusī, \emph{Ghayth al-Naf\textsuperscript{c} fī al-Qirāʾāt al-Sab\textsuperscript{c}}, introduction and passim.
\bibitem{madina} King Fahd Complex for the Printing of the Holy Quran, \emph{Muṣḥaf al-Madīna} (Ḥafṣ), closing note on its sources.
\bibitem{hamad} Gh. Q. al-Ḥamad, \emph{Taḥzīb al-Qurʾān fī al-Maṣādir wa-l-Maṣāḥif}.
\bibitem{masahif} Ibn Abī Dāwūd, \emph{Kitāb al-Maṣāḥif}.
\bibitem{bayan} Abū ʿAmr al-Dānī, \emph{al-Bayān fī ʿAdd Āy al-Qurʾān}.
\bibitem{muktafa} Abū ʿAmr al-Dānī, \emph{al-Muktafā fī al-Waqf wa-l-Ibtidāʾ}.
\bibitem{manar} A. al-Ushmūnī, \emph{Manār al-Hudā fī Bayān al-Waqf wa-l-Ibtidāʾ}.
\bibitem{tanzil} Tanzil Project, ``Quran text, Uthmani, version 1.1,'' \url{https://tanzil.net}.
\bibitem{digitalkhatt} A. Anane, ``DigitalKhatt,'' \url{https://digitalkhatt.org}.
\bibitem{tajemb} T. Eldeeb, ``tajweed-embeddings,'' \url{https://github.com/tarekeldeeb/tajweed-embeddings}.
\bibitem{nazm} T. Eldeeb, ``Quran-Nazm,'' \url{https://github.com/tarekeldeeb/quran-nazm}.
\bibitem{bellman} R. Bellman, ``On the approximation of curves by line segments using dynamic programming,'' \emph{Commun. ACM}, vol.~4, no.~6, p.~284, 1961.
\bibitem{repo} T. Eldeeb, ``Quran-Tajzee,'' \url{https://github.com/tarekeldeeb/quran-tajzee}.
\end{thebibliography}

\end{document}
